% /root/monodromy, 2026-11-19. Proved; independent root and reduction_descent audits passed.
\subsection{A level-three theta symmetrizer for the Cartier pencil}
\label{subsec:ppav-theta-Cartier-symmetrizer}
\status{Proved and checked. The parent and reduction/descent agent
independently audited the complete proof, opened the cited primary sources,
and reran the exact finite Heisenberg calculation. The relative Cartier
and theta-group identifications include the Frobenius twists explicitly.}

\begin{theorem}\label{thm:ppav-level3-Cartier-symmetric}
Let $(A,\Theta)$ be a principally polarized abelian variety of dimension
$g$ over an algebraically closed field $k$ of characteristic two, and let
$L=\mathcal O_A(3\Theta)$ and $W=H^0(A,L)$.
Choose a nonzero invariant top form $\omega$ and put
$C_{\rm abs}(\omega)=c^{1/2}\omega$, with $c\in k$.
There is a theta basis of $W$, indexed by $\mathbf F_3^g$, in which
the root-removed Cartier multiplication matrix $M(s)$ is symmetric
and has diagonal
\[
                         \operatorname{diag}M(s)=c\,s.
\]
If $A$ is ordinary, normalize $\omega$ by $C_{\rm abs}(\omega)=\omega$.
Then the diagonal coordinates are independent and $\det M(s)$ is a
nonzero geometrically squarefree polynomial of degree $3^g$.
In particular, for $g=2$ the degree-one full coherent Hasse divisor of
$|3\Theta|$ is geometrically reduced and each of its components has
generic linearized Cartier corank one.
\end{theorem}

\begin{proof}
\paragraph{Relative Cartier and its target.}
Let $F=F_{A/k}:A\to A^{(2)}$ be relative Frobenius. Put
$L^{(2)}=L\otimes_{k,F_k}k$ on the Frobenius twist and
$W^{(2)}=H^0(A^{(2)},L^{(2)})$. The canonical map
$\iota:W\to W^{(2)}$, $w\mapsto w^{(2)}$, is square-semilinear.
The relative top Cartier trace is $k$-linear and satisfies
\[
 C_{\rm rel}(\omega)=c\,\omega^{(2)},\qquad
 C_{\rm rel}(F^*a\,\eta)=a\,C_{\rm rel}(\eta).
\]
Since $L^2=F^*L^{(2)}$, multiplication followed by this trace defines
the $k$-bilinear symmetric map
\begin{equation}\label{eq:ppav-relative-beta}
 \beta:W\otimes W\longrightarrow W^{(2)},\qquad
 C_{\rm rel}(s t\omega)=\beta(s,t)\omega^{(2)}.
\end{equation}
In particular $\beta(s,s)=c\,s^{(2)}$.
These are the relative finite-duality and coefficient-trace formulas
of \cite[Lemma~1.3.6, Proposition~1.3.7, Theorem~1.3.8, pp.~26--28]
{BrionKumar2005}. The target is the Frobenius twist; no identification
of $A$ with $A^{(2)}$ is being made.

\paragraph{A finite theta subgroup and its two representations.}
The theta group of $L$ is a central extension of
$K(L)=A[3]\simeq\mathbf F_3^{2g}$ by $\mathbf G_m$ with nondegenerate
commutator pairing. Choose a Lagrangian decomposition and lifts of
its generators of order three. They, together with the central
$\mu_3$, generate a finite Heisenberg subgroup $H$ of order
$3^{2g+1}$. A primitive cube root $\zeta\in k$ identifies its
Schrodinger representation on a basis $e_a$, $a\in\mathbf F_3^g$, by
\[
 T_x e_a=e_{a+x},\qquad D_y e_a=\zeta^{y\cdot a}e_a,
 \qquad Z e_a=\zeta e_a.
\]
Here $T_x,D_y,Z$ generate $H$; their commutation convention merely
chooses the sign of the standard symplectic pairing.
The natural theta action on $W$ is this representation, denoted
$\rho$. The precise primary statements are
\cite[Corollary~8.20, Lemma~8.24, Theorem~8.28, description~8.29,
pp.~119, 121--123; Corollary~9.17, pp.~133--134]
{EdixhovenVanDerGeerMoonenAV}. They give the nondegenerate theta
group, its Heisenberg form and the unique weight-one representation
on $H^0(L)$; its dimension is $3^g$.

Relative Cartier is natural under translations. A theta lift on $L$
induces its tensor square on $L^2$; under $L^2=F^*L^{(2)}$, this
tensor square is the pullback of its Frobenius-twisted lift on
$L^{(2)}$. The invariant form is unchanged by translation.
Consequently, in the bases $e_a$ and $e_a^{(2)}$, one has the
fully typed equivariance identity
\begin{equation}\label{eq:ppav-Cartier-equivariance}
 \beta(\rho(h)s,\rho(h)t)=\rho(h)^{[2]}\beta(s,t)
                         \quad(h\in H).
\end{equation}
One can also verify this directly in a local frame: if a lift acts
by a line-bundle multiplier $u$, its action on a product has multiplier
$u^2$, and the Cartier projection formula moves its Frobenius
descendant to the target. On the base the translation descends from
$x$ to $F(x)$. Thus the target action is the coefficient-square
representation, not the original weight-one action.

On the displayed generators of the finite subgroup,
\[
                \rho(h)^{[2]}=\rho(h)^{-t}.
\]
For $T_x$ both sides are the same permutation matrix; for $D_y$ and
$Z$ it is the identity $\zeta^2=\zeta^{-1}$. Since both sides are
representations, it holds for every $h\in H$.
Therefore
\[
 R:W^{(2)}\longrightarrow W^*,\qquad e_a^{(2)}\longmapsto e_a^*,
\]
is an invertible $H$-equivariant linear map. This step uses the finite
central subgroup $\mu_3$. The corresponding statement for the whole
central $\mathbf G_m$ would be false and is not used.

\paragraph{Cyclic symmetry is a translation in level three.}
Define the trilinear tensor
\[
                  T(s,t,u)=\langle R\beta(s,t),u\rangle.
\]
It is $H$-invariant by \eqref{eq:ppav-Cartier-equivariance} and is
symmetric in its first two arguments by multiplication.
Write $T_{abc}=T(e_a,e_b,e_c)$. Invariance under every modulation
$D_y$ forces
\[
            T_{abc}=0\quad\text{unless }a+b+c=0.
\]
Invariance under $T_x$ gives
$T_{a+x,b+x,c+x}=T_{abc}$.
If $a+b+c=0$, translation by $x=b-a$ sends the ordered triple
$(a,b,c)$ to $(b,c,a)$, because $2=-1$ in the index group
$\mathbf F_3^g$. Thus $T$ is invariant under cyclic permutation.
Together with its first-two symmetry this proves full symmetry
under all permutations of its three arguments.

Let $M(s)$ represent $\beta(s,-)$ from the $e_a$ basis to the
$e_a^{(2)}$ basis. The matrix of $R$ in these chosen bases is the
identity. Full symmetry of $T$ gives $M(s)^t=M(s)$.
Furthermore $\beta(e_a,e_a)=c e_a^{(2)}$, so
$T_{aab}=c\delta_{ab}$. By full symmetry, the $a$th diagonal
entry of $M(s)$ is $c s_a$, proving the asserted formula.

For ordinary $A$, Cartier on its invariant top form is nonzero.
If $C_{\rm abs}(\omega)=u\omega$ initially, replacing $\omega$
by $u^2\omega$ normalizes it; hence $c=1$.
Lemma~\ref{lem:symmetric-independent-diagonal-squarefree} now
applies: the determinant has the full squarefree coordinate
monomial with coefficient one, while every nontrivial matching
term has an even positive exponent. It is nonzero and has no
repeated geometric irreducible factor.
At the generic point of a reduced determinant component the
matrix has corank one: corank at least two would make every
cofactor, and hence every first derivative of the determinant,
vanish. This finishes the proof.
\end{proof}

\begin{remark}[Coordinate changes and the finite check]
In an arbitrary original basis the conclusion is the existence of
a constant invertible matrix $R_0$ for which $R_0M(s)$ is symmetric
with independent diagonal coordinates. The matrix $R_0$ itself
need not be symmetric, because the change of basis on the target
is Frobenius-twisted. This is consistent with the explicit elliptic
product symmetrizers proved earlier.
The retained finite calculation
\path{agents/monodromy/code/level3_theta_tensor_checks.py}
checks the $g=2$ index action, all $243$ finite Heisenberg matrices,
and the full permutation symmetry of the five-dimensional space of
invariant first-two-symmetric tensors. It is a check of the finite
algebra, not a substitute for the geometric equivariance proof.
\end{remark}
